% Generated by scripts/materialize_unified_paper.py; do not edit.
\section{Reverse induction and perfect-sequence extraction}\label{low:sec:induction}

By \eqref{low:main:eq:7-7}--\eqref{low:main:eq:7-11}, after increasing the existential threshold \(C_0(K)\) and
restricting \(w\) to a sufficiently small right-neighborhood of
\(\omega_0\), we have
\(\widehat B_w(C)>0\), \(E_w(C)>0\), and \(\mu_w>0\).  We fix this range
throughout the section; in particular, every monotonicity assertion below
has its sign justified before it is used.

Define the new accumulated deterministic deficit
\[
 \mathcal N_{C,w,d}(r)=\frac1d\left[
 2\widehat B_wA_0\binom r4+E_wK_r\right].
\tag{L6.1}\label{low:main:eq:6-1}
\]
At an exposure step with \(k=r-i\) future vertices, the same-history
calculation of \cref{low:sec:ledger}, followed by the pathwise boundary estimate
\eqref{low:main:eq:4-6}, proves the one-step kernel bound \eqref{low:main:eq:2-37} with
\[
 \lambda_i=\mathsf d_{k,w,d}
 +\widehat B_w(C)\frac{S_i}{d}
 -\sum_{j>i}\frac{e_{ij}^2}{2\mu_wd}
 -\frac{d^{-1/5}}{\sqrt d}\sum_{j>i}T_{ij}.
\tag{L6.2}\label{low:main:eq:6-2}
\]
Every coefficient is deterministic and fixed before column \(i\) is
sampled.  Equations \eqref{low:main:eq:2-33} and \eqref{low:main:eq:2-40} retain every future coefficient
\(T_{jq}\), while \eqref{low:main:eq:4-6} supplies each new boundary edge \(ij\) with its full
coefficient \(T_{ij}\).  Thus \cref{low:lem:weighted-reverse} applies without an
exchangeability assumption.

Apply \eqref{low:main:eq:2-38} at \(s=0\), average over \(\mathcal F_0\), and then use
\cref{low:lem:multiplicity}.  Since \(\mathcal P_0^T=0\), this gives
\[
P^*_{R,r}\le\mathcal H_C(r)
\exp\!\left[-\mathcal D_{C,w,d}(r)
-\mathcal N_{C,w,d}(r)+J_{r,d}\right],
\tag{L6.3}\label{low:main:eq:6-3}
\]
where
\[
0\le J_{r,d}\le
\frac{m_0+\theta/D}{1-\rho}d^{-1/5}\binom r2=o_C(r^2).
\tag{L6.4}\label{low:main:eq:6-4}
\]
Here \(J_{r,d}\) denotes the displayed deterministic upper envelope, so
\(J_{q,d}\le J_{\ell,d}\) for \(q\le\ell\).  The base cases
\(r=0,1,2\) follow from the empty-event convention and the same one-step
identity; no asymptotic estimate is needed there.

\subsection{Uniform restoration after deleting vertices}

Let extraction retain \(q=\ell-u\) vertices.  Relabel them in retained order
and rebuild the deterministic array for target size \(q\).  Since
\(q/\sqrt d\le\ell/\sqrt d\), all row-sum, Hessian and H\"older feasibility
inequalities only improve.

First, the bracket in \eqref{low:main:eq:2-14} lies in \([0,1]\), and \(m_w\le m_0\).
Consequently
\begin{align}
0\le{}&\mathcal D_{C,w,d}(\ell)
-\mathcal D_{C,w,d}(\ell-u)\notag\\
\le{}&\frac{m_0^4}{d}
\sum_{k=\ell-u}^{\ell-1}k(k-1)^2
\le\frac{m_0^4}{D^2}\,\ell u.
\tag{L6.5}\label{low:main:eq:6-5}
\end{align}
For sufficiently large \(C\),
\(m_0^4/D^2=(1+o(1))/(8\log C)<1\), so the right side is at most
\(\ell u\).

Next,
\[
 \binom\ell4-\binom{\ell-u}4\le u\binom\ell3,\qquad
 \binom\ell3-\binom{\ell-u}3\le u\binom\ell2.
\tag{L6.6}\label{low:main:eq:6-6}
\]
Using \(K_r=\binom r3+2\binom r4\), \eqref{low:main:eq:6-6} yields
\begin{align}
0\le{}&\mathcal N_{C,w,d}(\ell)
-\mathcal N_{C,w,d}(\ell-u)\notag\\
\le{}&\frac1d\left[
2\widehat B_wA_0u\binom\ell3
+E_w\left(u\binom\ell2+2u\binom\ell3\right)\right]\notag\\
\le{}&\left[
\frac{\widehat B_wA_0+E_w}{3D^2}
+\frac{E_w}{2D^2\ell}\right]\ell u.
\tag{L6.7}\label{low:main:eq:6-7}
\end{align}
The last bracket is \(O(1/\log C)+O_C(\ell^{-1})<1\) for sufficiently large
\(C\) and then sufficiently large \(\ell\), so
\eqref{low:main:eq:6-7} also costs at most \(\ell u\).

Finally, the ratio between the source envelopes at sizes \(\ell-u\) and
\(\ell\) costs at most
\[
 \left[\log(1/p)+1\right]\ell u.
\tag{L6.8}\label{low:main:eq:6-8}
\]
Indeed, \(\binom\ell2-\binom{\ell-u}2\le\ell u\) pays the missing
\(p\)-factors, and
\(\binom\ell3-\binom{\ell-u}3\le\ell^2u/2\), together with
\(\vartheta_Ca^3/(p^3D)\le2\) for all sufficiently large \(C\), pays the
cubic term.  The positive \(d^{-1/5}\binom r2\) and \(r\) terms in
\(\mathcal H_C(r)\), as well as \eqref{low:main:eq:6-4}, are monotone in \(r\) and require no
restoration charge.

\subsection{The binomial subset sum}

Apply the pinned extraction inequality \eqref{low:main:eq:2-41}, then \eqref{low:main:eq:6-3} at
\(q=\ell-u\), and finally \eqref{low:main:eq:6-5}--\eqref{low:main:eq:6-8}.  For every specified retained set
\(I\) and every \(1\le u\le\ell\),
\[
\begin{aligned}
p_I
&\le P^*_{R,\ell-u}\left(\frac p{10}\right)^{10\ell u}\\
&\le\mathcal H_C(\ell)
\exp\!\left[-\mathcal D_{C,w,d}(\ell)
-\mathcal N_{C,w,d}(\ell)+J_{\ell,d}\right]
e^{-\eta_C\ell u},
\end{aligned}
\tag{L6.9}\label{low:main:eq:6-9}
\]
where the explicit residual exponent
\[
 \eta_C=10\log(10/p)-\log(1/p)-3>0
\tag{L6.10}\label{low:main:eq:6-10}
\]
pays, respectively, the source-envelope, old-CGF, and new-residual
restoration charges.  Summing over all possible deleted sets is now an exact
binomial calculation:
\[
 \sum_{u=1}^{\ell}\binom\ell u e^{-\eta_C\ell u}
 =(1+e^{-\eta_C\ell})^\ell-1=o(1).
\tag{L6.11}\label{low:main:eq:6-11}
\]
Including the \(u=0\) perfect term proves the unconditional red probability
bound
\[
P_{R,\ell}\le\mathcal H_C(\ell)
\exp\!\left[-\mathcal D_{C,w,d}(\ell)
-\mathcal N_{C,w,d}(\ell)+J_{\ell,d}\right](1+o(1)).
\tag{L6.12}\label{low:main:eq:6-12}
\]
This proves extraction of the complete new exponent, rather than merely
asserting that a source failure factor is large enough.
